\ficha{Fonseca e Mollo (2012), Metalistas x papelistas}{fonsecamollo2012}

\begin{identificacao}
FONSECA, Pedro Cezar Dutra; MOLLO, Maria de Lourdes Rollemberg. Metalistas x
papelistas: origens teóricas e antecedentes do debate entre monetaristas e
desenvolvimentistas. \emph{Nova Economia}, Belo Horizonte, v. 22, n. 2, p.
203-233, maio-ago. 2012.

Artigo de história do pensamento econômico, 31 páginas, em português. Lido
integralmente na versão publicada.
\end{identificacao}

\bloco{Problema e tese}
O artigo pergunta de onde vêm, teoricamente, as categorias que organizaram o
debate monetário brasileiro da segunda metade do século XIX, e o que elas
produziram depois. A resposta tem duas pernas. Metalismo e papelismo são
reflexo tardio das duas controvérsias monetárias inglesas, bulionistas contra
antibulionistas e Currency School contra Banking School, cujo eixo comum é a
neutralidade ou não neutralidade da moeda. E o papelismo, ao adaptar a
herança heterodoxa a uma economia agroexportadora e ao optar pelo crescimento
contra a estabilidade, é uma das fontes do desenvolvimentismo brasileiro,
tese que os autores testam na trajetória de Getúlio Vargas.

\bloco{Argumento}
\begin{enumerate}[leftmargin=*,nosep]
\item Primeira controvérsia inglesa, 1797 a 1825. Os bulionistas, com Ricardo
à frente, e ainda Wheatley e Thornton, partem da Teoria Quantitativa da Moeda
e leem o prêmio do \emph{bullion} sobre o preço de cunhagem como sinal de
excesso de emissão, o que faz da conversibilidade-ouro a regra monetária de
controle de emissão e de preços (p. 205). Os antibulionistas, Bosanquet,
Torrens e Boase, atribuem o prêmio e o déficit externo a causas externas,
remessas militares e quebras de safra, à variação da velocidade de circulação
e à \emph{real bills doctrine}, e temem a contração porque percebem seus
efeitos sobre a produção (p. 206).

\item O corte analítico do artigo é a neutralidade. Nos bulionistas, a
neutralidade é o que permite não esperar efeito duradouro da moeda sobre a
produção real, e daí a prioridade ao controle de preços; nos antibulionistas,
o receio do impacto sobre emprego e produção é a própria percepção da não
neutralidade (p. 207). O mais heterodoxo do período é Thomas Attwood, da
Birmingham School, contrário ao automatismo do padrão-ouro por suas
consequências sobre crescimento e emprego (p. 213).

\item Segunda controvérsia, 1825 a 1875, com o que os autores chamam de
ortodoxização. Currency School e Banking School aceitam a conversibilidade-ouro
como regra necessária, e o ponto de divergência passa a ser a necessidade de
controles de curto prazo (p. 209). A Currency School propõe o \emph{currency
principle}, a Palmer Rule e a lei bancária de 1844, para fazer o sistema misto
funcionar como puramente metálico. A Banking School argumenta que depósitos e
letras também são moeda, que a quantidade de moeda depende do ritmo dos
negócios e não o inverso, e que a lei do refluxo drena qualquer excesso, de
modo que os controles são inviáveis e indesejáveis (p. 210-211). Os autores
observam que a posição da Banking School é mais enfática na endogeneidade da
moeda do que na sua não neutralidade (p. 212).

\item A dobradiça com o Brasil vem de Dean. A dominação ortodoxa durou tanto
na Inglaterra por causa da posição hegemônica do país, que lhe dava superávits
confortáveis, e por isso as ``limitações automáticas sobre a política de
crédito interno só raramente foram muito exigidas'' (p. 213). Não é o caso do
Brasil, que enfrentava, como outros países menos desenvolvidos, os sacrifícios
das exigências de reserva. A dificuldade prática de sustentar o padrão-ouro na
periferia é o que explicita a não neutralidade da moeda por outra via.

\item O debate brasileiro é singular por duas razões: a economia
agroexportadora e o deslocamento do foco, da velha polaridade entre
liberalismo e intervencionismo para o \emph{modus faciendi} da política
econômica, seus objetivos de curto prazo e seus instrumentos (p. 215). O
epicentro é a conversibilidade.

\chave{Os metalistas defendiam ferrenhamente o padrão-ouro e a
conversibilidade da moeda; para tanto, encontravam respaldo na teoria
econômica convencional e na política do país hegemônico, a Grã-Bretanha. Já os
papelistas, ante a ausência de um corpo teórico de mesma envergadura [...]
recorriam à razão prática: a experiência, e não uma teoria, demonstrava qual o
melhor caminho a seguir}{p. 215}

\item A filiação é afirmada com prudência. Metalistas e papelistas aparecem na
segunda metade do século XIX, com defasagem em relação ao debate inglês, e os
autores dizem ser ``razoável supor'' que este os tenha inspirado, dada a
hegemonia britânica também intelectual (p. 217). O papelismo se divide em dois
grupos, análogos aos meios-termos de Schumpeter: moderados, que não negam a
conversibilidade como regra mas pedem relaxamento nas crises e nas safras ou
lastro proporcional flexível, entre eles Souza Franco, Mauá, os viscondes do
Cruzeiro e de Ouro Preto, João Alfredo e Lafayette; e radicais, como Rui
Barbosa, que negam qualquer regra de conversibilidade e defendem pluralidade
de bancos emissores (p. 217).

\item A oposição se resume numa inversão de hierarquia entre variáveis. Para
os metalistas, causa interna, política monetária subordinada à cambial, e a
taxa de câmbio como variável prioritária (p. 218). Para os papelistas, a
atenção maior deveria estar na taxa de juros, e não na taxa de câmbio, com a
política cambial subordinada à monetária e esta às necessidades da produção
(p. 220). Mauá formula o ``requisito da elasticidade'', a oferta de moeda
flexível a ponto de não interferir negativamente nas atividades produtivas
(p. 218). Rui Barbosa nega que a circulação metálica firme o câmbio e afirma o
inverso, que a prosperidade produz o câmbio favorável e este permite a
circulação conversível (p. 220-221), e conclui que o ``barômetro das
exagerações do meio circulante não é a taxa de câmbio'', é a taxa de juros
(p. 221).
\end{enumerate}

\chave{O papelismo, desta forma, representou uma precoce heterodoxia ao
redefinir quem era o cão e quem era a cauda [...] e, com isso, priorizar o
investimento sobre a poupança, a taxa de juros sobre a taxa de câmbio e o
crescimento sobre a estabilidade}{p. 222}

\bloco{Conceitos e definições operacionais}
Metalismo: defesa do padrão-ouro e da conversibilidade como regra, com
prioridade à estabilidade de preços e ao câmbio, apoiada na TQM e na economia
clássica (p. 215, p. 218). Papelismo: prioridade ao nível de atividade, com a
pergunta operacional sobre qual oferta monetária é mais condizente com o ânimo
dos negócios, e apelo à experiência em lugar de doutrina (p. 218). Neutralidade
da moeda: critério que separa os campos, definido pela ausência ou presença de
efeito duradouro da moeda e do crédito sobre a produção real (p. 207, p. 222).
Endogeneidade: quantidade de moeda determinada pelo ritmo dos negócios, não o
inverso (p. 211). Lei do refluxo: empréstimos saldados retornam ao sistema
bancário, e o excesso de notas é convertido e drenado, o que dispensa controle
quantitativo (p. 210-211). Requisito da elasticidade: oferta de moeda flexível
o bastante para não travar a produção (p. 218). Intervencionismo
pró-crescimento: uso de política monetária, cambial e fiscal para acelerar o
crescimento, um dos três itens do núcleo duro do desenvolvimentismo (p. 224).

\bloco{Evidência e método}
Duas camadas distintas. A parte inglesa é reconstrução de história do
pensamento a partir de fontes secundárias e de compilações, Viner (1937),
O'Brien (2004), Schumpeter (1954), Dean (1980), Schwartz (1989) e Mollo
(1994), com citações de Ricardo, Boase, Attwood, Overstone, Tooke e Marx
tomadas quase sempre de segunda mão. A parte brasileira do século XIX também é
secundária, apoiada em Franco (1983), Gremaud (1997), Neuhaus (1975), Müller
(2004) e Prado (2003), com falas de Afonso Celso, Lafayette e Rui Barbosa
citadas por edições e compilações. Só a última seção é pesquisa em fonte
primária, e o objeto é Vargas: discursos e mensagens de 1906, 1919, 1927 e
1928 no \emph{Correio do Povo} e nos anais gaúchos, mais um texto de 1950. Não
há série quantitativa, teste nem contrafactual.

\bloco{Agentes e vertentes}
Inglaterra, lado ortodoxo: Ricardo, Wheatley, Thornton como bulionista
moderado, e depois Torrens e Lord Overstone na Currency School. Lado
heterodoxo: Bosanquet e Boase entre os antibulionistas, Tooke e Fullarton na
Banking School, Attwood e a Birmingham School como polo mais radical, e Marx
lido como mais heterodoxo do que a Banking School. Brasil, metalistas:
Francisco Belizário, Torres Homem e Joaquim Murtinho, ministro da Fazenda de
Campos Sales. Papelistas moderados: Souza Franco, Mauá, os viscondes do
Cruzeiro e de Ouro Preto, João Alfredo e Lafayette. Papelista radical: Rui
Barbosa. Herdeiros: Vargas, e a leitura do desenvolvimentismo em Bielschowsky,
Furtado, Prebisch e Caio Prado Jr.

\bloco{Críticas e limites}
O recorte para no fim do Império e nos primeiros anos da República, e depois
salta para os anos 1920 e 1930. A Caixa de Conversão, o Convênio de Taubaté e
o período de 1906 a 1914 não são tratados: a única passagem de 1906 é um
discurso de estudante de Vargas, alheio ao debate monetário. A filiação entre
as escolas inglesas e as correntes brasileiras é apresentada como suposição
razoável a partir da hegemonia britânica (p. 217), sem evidência de leitura,
tradução ou circulação de textos. A base social de cada corrente é conjectura
assumida, na ausência de estudos empíricos conclusivos (p. 216). Os próprios
autores relativizam a dicotomia duas vezes, ao lembrar que a Banking School
defendia a conversibilidade-ouro e ao acolher de Gremaud (1997) a existência
de formas intermediárias dentro de um mesmo regime cambial (p. 219-220). A
prova do elo com o desenvolvimentismo repousa numa única trajetória, num único
estado, e o critério de semelhança de argumentos não exclui convergência
independente.

\begin{dialogo}
\item Sobrevida do vocabulário. Procurar nos quatro jornais as palavras
``metalista'' e ``papelista'' entre 1906 e 1914, e verificar quem as usa e
contra quem. Fonseca e Mollo situam a controvérsia no Império e na Primeira
República inicial. Se os termos aparecerem em 1906 apenas como insulto ou
memória do Encilhamento, isso sustenta a hipótese do capítulo de que a
clivagem já não organiza os campos e foi substituída pelo eixo valorização
contra emissão.

\item Requisito da elasticidade. Procurar a defesa da Caixa formulada como
necessidade de circulação flexível na safra e na entressafra, com vocabulário
de época do tipo ``necessidades do comércio'', ``ânimo dos negócios'' e
``meio circulante''. Testa se o argumento papelista migra intacto para o campo
pró-Caixa, ou se a defesa da Caixa se faz por argumento novo, a estabilização
da taxa a um nível dado.

\item Hierarquia das variáveis. Verificar se algum publicista de 1906 a 1914
sustenta, como Rui Barbosa em 1891, que a variável-chave é a taxa de juros e
não a de câmbio. O debate da Caixa gira sobre o nível da taxa em pence, o que
é o oposto da inversão papelista. Se a taxa de juros estiver ausente da
argumentação de Campista, Bulhões, Serzedello Corrêa e Alcindo Guanabara, o
papelismo não sobreviveu como quadro analítico, apenas como rótulo.

\item Experiência contra doutrina. Procurar a fórmula que opõe a prática às
teorias importadas, e o argumento dos ``países novos'' contra a universalidade
das leis econômicas, que Fonseca e Mollo extraem de Mauá (p. 219). Testa se a
defesa da Caixa se apresenta como pragmatismo sem teoria, como o papelismo, ou
se reivindica respaldo doutrinário, o que indicaria deslocamento do repertório.

\item Rui Barbosa em 1906. O artigo o classifica como papelista radical em
1891. Levantar suas manifestações sobre a Caixa nos jornais e nos Documentos
Parlamentares. Uma posição de Rui contrária à emissão contra ouro, ou
indiferente a ela, seria evidência direta de que o eixo se reconfigurou entre
1891 e 1906.
\end{dialogo}

\usonocapitulo{Parte III, como texto de abertura. É a referência que autoriza
tratar metalismo e papelismo como importação das controvérsias inglesas e não
como invenção local, e que fornece o critério com que os autores organizam as
duas margens, a neutralidade da moeda. Serve a três afirmações do capítulo.
Primeira, que a ortodoxia monetária chega ao Brasil pela via da autoridade
teórica e do exemplo do país hegemônico, o que liga a Parte III às Partes I e
II. Segunda, que a assimetria periférica muda o custo da regra, com o
argumento de Dean sobre a folga que os superávits davam à Inglaterra (p. 213).
Terceira, que a opção papelista era por crescimento e pela inversão da
hierarquia entre juros e câmbio, o que dá a régua para medir o que resta dessa
posição em 1906. O artigo não trata da Caixa de Conversão nem do período da
dissertação, e nenhuma afirmação sobre 1906 a 1914 pode ser apoiada nele. Sua
função é fornecer a genealogia das categorias e, por contraste, sustentar a
tese do capítulo de que a clivagem do Encilhamento não é a de 1906.}
